\documentclass[11pt]{article}
\usepackage[margin=1in]{geometry}
\usepackage{amsmath,amssymb}
\usepackage{parskip}
\usepackage{xcolor}
\usepackage{booktabs}
\usepackage{longtable}
\usepackage[hidelinks]{hyperref}
\usepackage{fancyvrb}
\newcommand{\code}[1]{\texttt{#1}}
\definecolor{warn}{rgb}{0.65,0.10,0.10}
\definecolor{ok}{rgb}{0.10,0.45,0.15}
\newcommand{\stopsign}[1]{{\color{warn}\textbf{#1}}}
\newcommand{\good}[1]{{\color{ok}\textbf{#1}}}

\title{Non-stationary conditional extremes for BC streamflow:\\
presenting the analysis, and what is inside the code}
\author{Alex Lin \quad (supervisor: Natalia Nolde)}
\date{August 2026 \quad (version 4)}

\begin{document}
\maketitle

\noindent
This document exists for two jobs. \textbf{Job 1: present the analysis.}
Part~I walks the report in the order I would show it; for every figure it
gives the axes, the theory behind the picture (with the derivation where one
exists), the exact mechanism that produced it, what our own numbers say, and
the follow-up questions that figure invites, with answers. \textbf{Job 2:
prove the code is not a black box.} Part~II opens each algorithm: what it
computes, the pseudocode, why that is the right computation, which part is a
published package and which part is ours, and how each is verified. Part~III
is a question-and-answer drill. Part~IV is the fix list.

Pointers of the form ``v8~\S$x$, eq.~($n$), p.~$y$'' give the section,
equation and page of \emph{my current understanding of the model}, version~8.
Every number quoted is from the knit of \code{code\_report\_v4.Rmd} on
\code{daily\_sf\_bf\_rcp45\_only\_1945\_2099\_9avg.csv}: R 4.5.2, qgam 2.0.0,
evgam 1.0.1, mgcv 1.9-3, texmex 2.4.9, extRemes 2.2-1, quantreg 6.1; axis
1969-01-01 to 2099-12-31, $n = 47{,}847$ days; conditioning pair BRIDG given
NTHMB (Case~4), anchor BARRM given BARRS (Case~1).

\tableofcontents
\newpage

\section{Status: what is ready to show, and what is not}

Read this page before the meeting. The analysis has two halves, and they are
in different states.

\begin{longtable}{@{}p{4.3cm}p{2.1cm}p{8.2cm}@{}}
\toprule
\textbf{Stage} & \textbf{State} & \textbf{Basis} \\
\midrule
\endhead
Data cleaning gate & \good{present} & 0 duplicate rows, 0 missing dates, 0 NA,
no non-positive baseflow; streamflow effectively tie-free (most repeated value
occurs twice; 0.0\% of days on values repeated $\ge 10$ times) \\
Seasonal threshold & \good{present} & Threshold tracks the annual cycle in
every window; trend slopes significant at all three series ($z = 11.8$, $7.0$,
$17.7$) \\
Local 5\% calibration & \good{present, with the upgrade} & Additive design
misses the first era (SON 0.149, JJA 0.123); the tensor
\code{te(doy,\,tt)} repairs it (SON 0.039, JJA 0.039) \\
GP tail & \good{present} & 2236--2299 exceedances per series; scale trend
significant everywhere ($p<10^{-8}$) \\
PIT / Laplace gate & \good{present} & Histograms flat to within the noise the
sample allows (Section~\ref{sec:pit} quantifies ``flat''), sliced QQ on the
line in both seasons and both decades, no point mass \\
Extremal index, return levels, machinery checks & \good{present} & extRemes
$\equiv$ project to three decimals; Route A vs B agree to $\le 1\%$ \\
\midrule
\stopsign{Dependence stage} & \stopsign{do not present as a finding} &
The constrained optimiser did not move off its starting value: all three
model sizes returned the identical objective 2698.6, the trend test returned
exactly 0.00, $\hat a(t)$ came out exactly flat. Section~\ref{sec:defect}
diagnoses it and gives the repair \\
\stopsign{Conditional functionals} & \stopsign{report as provisional} &
They are computed from that stuck dependence fit, so the flat
$P_{co}\approx 0.35$ is a property of the defect, not of the rivers \\
\bottomrule
\end{longtable}

\noindent
The honest sentence for the meeting: \emph{the marginal stage passes every
check you asked for; the dependence stage has a numerical defect I have
located and can state precisely, and its results are on hold until it is
fixed.} That is a stronger position than presenting a number that cannot be
defended --- and Section~\ref{sec:defect} shows the defect is already
half-solved, because two independent pieces of evidence agree on what the
answer should look like.

\section*{Part I --- The walkthrough}
\addcontentsline{toc}{section}{Part I --- The walkthrough}

\subsection{Station 1: the data gate}
\label{sec:datagate}

\textbf{What is shown.} Printed counters (duplicates, gaps, missing values,
non-positive baseflow), then three forensic instruments per series: the ten
most repeated exact values with the months and years they fall in; runs of
identical consecutive days; and the audit of days at or near the series
minimum.

\textbf{Why ties are fatal if present --- the exact mechanism.} The whole
pipeline stands on the probability integral transform
(Section~\ref{sec:pitthm}), and that theorem requires a \emph{continuous}
distribution. Suppose instead some value $x_0$ occurs on many days, so
$\Pr(X = x_0) = p > 0$. Then the transformed variable $U = F(X)$ inherits an
atom: a mass of size $p$ sits at the single point $F(x_0)$, no matter how well
the model fits. On the histogram that is a spike; on the QQ plot it is a flat
run of identical sample values ending in a jump; on the Laplace scatter it is
a horizontal or vertical line of points. So the forensics are not tidiness ---
they are a precondition of the method, and each instrument targets one way
mass can enter: the repeated-value table finds atoms wherever they occur; the
run table finds \emph{plateaus} (one value frozen across consecutive days,
which is an instrument or post-processing signature, since rivers vary daily
--- and which also inflates serial dependence); the floor audit finds a hard
minimum (detection limit or frozen gauge), the single most common atom in
flow data.

\textbf{How each table is computed.} The repeated-value table is a frequency
count of exact values, keeping the top ten, with each value's day count, year
range, number of distinct years, and the months it occurs in --- the last two
columns answer the professor's question ``which values are these, and is it
only one year or many.'' The run table applies run-length encoding to the
daily series and keeps every run of $\ge 3$ equal consecutive values with its
start date, length and value, plus a month table of where runs live. The
floor audit counts days exactly at the minimum and within 1\% of it.

\textbf{Our numbers.} Streamflow, both stations: the most repeated value
occurs exactly \emph{twice} in 47{,}847 days; 0.0\% of days sit on values
repeated ten or more times; zero runs of three identical days; the minimum is
carried by a single day. Baseflow: 3.0\% of days (1447) carry a value repeated
$\ge 10$ times --- the repeated values are spread over 10--14 different years
each and concentrate in low-flow months --- and there are exactly four short
plateaus (lengths 3, 3, 3, 4; one in February, three in March). File-level:
0 duplicated (station, date) rows, 0 missing dates, 0 missing values, 0 days
of $\mathrm{bf} \le 0$.

\textbf{What to say.} ``The cleaning items are closed: no duplicated records,
no gaps, no missing values, no non-positive baseflow, and the streamflow
series carry no repeated values and no frozen stretches. The remaining tie
mass is confined to baseflow --- three percent of days, spread across years,
in the low-flow range --- which is a precision signature of the baseflow
product, matters only for Structure~B, and is documented rather than
hidden.''

\textbf{Question: ``how do you know the winter data are not corrupted?''}
Three pieces of evidence, all printed: no winter value repeats (streamflow),
no frozen stretches exist anywhere in streamflow, and the winter window
figures show the smooth recessions hydrology produces rather than flat
segments or discrete steps.

\textbf{Question: ``would you exclude the baseflow ties?''} No --- they are
endemic (every era, not one bad year), so exclusion would delete a property
of the data source rather than an error. The honest treatments are to note
the granularity in the paper and, if Structure~B tail diagnostics ever show
an atom-induced artifact, to jitter within the reported precision as a
sensitivity check.

\subsection{Station 2: the seasonal moving threshold}
\label{sec:threshold}

\textbf{What is shown.} Twenty separate panels --- both stations, streamflow
and baseflow, five randomly drawn two-year windows each. Grey dots: daily
values. Red: the working threshold $u(t; 0.95)$. Blue: the moving median
$u(t; 0.5)$. Each title carries that window's own share of days above red.

\textbf{The principle: a threshold is a probability, not a level.} The object
is the conditional quantile curve defined by
$\Pr\{X_t \le u(t;\tau)\} = \tau$ \emph{at every date} (v8~\S2.1, eq.~(3),
p.~3). On a snowmelt river a fixed horizontal level is a once-a-decade event
in February and a routine event in July, so its exceedances would be ``July
days,'' not ``extreme days,'' and every downstream stage would inherit that
confusion. Fixing the probability and letting the level move makes
``exceedance'' mean the same thing --- rarity $1-\tau$ --- on every date. That
is Lee and Lee's argument, and the local exceedance checks of
Station~\ref{sec:local5} verify the property rather than assume it.

\textbf{Why quantile regression finds a quantile --- the derivation.} The
fitted curve minimises the pinball (check) loss
$\rho_\tau(r) = \tau r$ for $r \ge 0$ and $(\tau-1) r$ for $r < 0$, applied to
residuals $r = x - u$. Differentiate the expected loss in $u$:
\[
\frac{d}{du}\,\mathbb{E}\,\rho_\tau(X-u)
= -\tau \Pr(X > u) + (1-\tau)\Pr(X < u)
= F(u) - \tau ,
\]
which is zero exactly when $F(u) = \tau$. So minimising the pinball loss
\emph{is} estimating the $\tau$-quantile --- not a heuristic, an identity.
Everything else (splines, penalties, qgam) is machinery for letting $u$
depend on $t$ smoothly while this identity does the statistical work.

\textbf{What the curve is made of.} Two components in one formula
(v8~\S4.1--4.2, pp.~5--7). A \emph{cyclic} day-of-year spline
(\code{bs = "cc"}): basis functions periodic with period one year, built so
the curve and its derivatives match across the 31~December / 1~January join
--- one seasonal shape, learned from all 131 years at once. Plus a slow
\emph{trend} spline in the running date: cubic bumps with knots spread over
the century, each coefficient shaped by the several years under its own bump.
The tensor upgrade of Station~\ref{sec:local5} replaces ``seasonal shape
$+$ trend'' by a smooth \emph{surface} in (day-of-year, year), letting the
shape itself evolve.

\textbf{Why not fit each calendar day separately.} Day-of-year alone has 366
covariate values and deletes the year --- June~14 of 1970 and of 2095 would be
forced to share one distribution, and change over the record (the object of
study) could not appear. The running axis keeps the year but has no repeats;
smoothness is what converts neighbouring days into information about a given
date. The two-component basis takes the best of both: repetition across years
for the cycle, neighbours along the axis for the trend.

\textbf{The trend test behind the slope numbers (move one).} The pinball loss
is a loss, not a likelihood, so there is no likelihood-ratio test here.
Instead: fit the rung with a linear trend term (harmonics included in
\emph{both} candidates, so the slope is trend, not season), then resample
whole years with replacement 30 times, refit, and use the spread of the 30
slopes as the standard error. Whole years, because flood days cluster ---
resampling days would understate the error. Our numbers: NTHMB slope $61.8$,
SE $5.2$, $z = 11.8$; BRIDG $7.12/1.02$, $z=7.0$; BRIDG baseflow
$1.77/0.10$, $z = 17.7$. All three high quantiles are rising, decisively.

\textbf{What to say.} ``The threshold tracks the seasonal cycle within each
year --- these are the one-to-two-year windows you asked for --- and each
window's own exceedance share is printed in its title. The trend in the high
quantile is significant at every series.''

\textbf{Question: ``what should the share in a title be?''} 0.05, with room
for noise: a two-year window holds about 730 days, so even independent days
give a standard deviation of $\sqrt{0.05 \times 0.95 / 730} \approx 0.008$,
and clustering roughly doubles it --- so anything in about $0.03$--$0.07$ is
consistent with perfect calibration, and a title far outside that range
localises a miscalibrated stretch.

\subsection{Station 3: the local 5\% table --- and why the tensor design wins}
\label{sec:local5}

\textbf{What is shown.} The share of days above $u(t;0.95)$ by season (rows)
and 22-year era (columns) --- 24 cells, target 0.05 in each --- under the
additive threshold; then the same table under the tensor threshold.

\textbf{Why this is the decisive check.} Pooled calibration is guaranteed
almost by construction (a quantile fit places 5\% of \emph{all} days above
itself); the professor's demand was ``roughly 5\% \emph{locally}, not just
overall,'' i.e.\ that the table be flat. And the check is \emph{out of
model}: no fitting criterion optimises these cells, so improvement in them is
evidence, not bookkeeping.

\textbf{How much a cell may wobble.} A cell holds roughly
$22 \times 90 \approx 2000$ days. Independent days would give the share a
standard deviation of $\sqrt{0.05 \times 0.95/2000} \approx 0.005$;
clustering inflates it to roughly $0.01$--$0.015$, so cells within about
$\pm 0.03$ of the target are noise. Cells at $0.123$ and $0.149$ are not
noise --- they are seven to ten standard deviations out.

\textbf{Our numbers.} Additive design, first era [1969,1991): SON $0.149$,
JJA $0.123$, MAM $0.074$, DJF $0.010$ --- the seasonal shape fitted as the
131-year average does not match the first decades, so the threshold sits too
low in summer--autumn and too high in winter there; later eras are close to
target. Tensor design \code{te(doy, tt)}, same first era: SON $0.039$, JJA
$0.039$, MAM $0.043$, DJF $0.012$; all 24 cells lie in $0.012$--$0.053$.

\textbf{The concept behind the tensor.} Additive
$f(\mathrm{doy}) + g(\mathrm{tt})$ means \emph{one} seasonal profile, the
same in 1970 and 2095, shifted up or down by the trend. The tensor product
builds a smooth surface from products of the two bases, so the freshet can
arrive earlier, sharpen or flatten across the century --- which on a warming
snowmelt river is precisely the expected physics. The additive model is a
special case of the tensor, so nothing is lost by the generalisation, and the
smoothing penalty (chosen by marginal likelihood) shrinks back toward the
simpler shape unless the data pay for the extra freedom.

\textbf{What to say.} ``The additive threshold fails your local-5\% test in
the first era; letting the seasonal shape itself drift --- a one-line change
of formula --- repairs every cell. I propose adopting the tensor design.''

\subsection{Station 4: the GP tail}
\label{sec:gp}

\textbf{What is shown.} A table (route, exceedance count, shape $\hat\xi$,
trend test), the fitted scale $\sigma(t)$ averaged by era, and two boxplot
panels: raw excesses by month, and excesses divided by the fitted
$\sigma(t)$, by month.

\textbf{Why generalized Pareto and nothing else.} The Pickands--Balkema--de
Haan theorem: for essentially any underlying distribution for which maxima
have a limit law, the distribution of the \emph{excess over a high
threshold}, $X-u$ given $X>u$, converges as the threshold rises to a
generalized Pareto distribution,
\[
\Pr\{X_t > x \mid X_t > u(t), t\}
= \Bigl(1 + \xi\,\frac{x-u(t)}{\sigma(t)}\Bigr)_{\!+}^{-1/\xi}
\qquad \text{(v8~\S2.2, eq.~(4), p.~3).}
\]
That is the licence to model only the tail and to extrapolate beyond the
data: the GP is not one convenient choice among many, it is the only
continuous limit available. The threshold is our moving $u(t;0.95)$, so the
theorem is applied date by date.

\textbf{Reading the two parameters.} $\sigma(t)$ is the typical overshoot
size at date $t$ --- a trend in it says floods are growing. $\xi$ is the
tail-weight: $\xi > 0$ heavy-tailed with no upper limit (Pareto-like);
$\xi = 0$ exponential; $\xi < 0$ a \emph{finite} upper end point at
$u + \sigma/|\xi|$. Letting $\xi$ move with $t$ would claim the character of
the tail is changing --- a far stronger claim than growth --- and $\xi$ is the
hardest parameter to estimate, so it is held constant per series (the
Faulkner et al.\ decision).

\textbf{Our numbers.} Exceedances 2264 / 2236 / 2299. Shapes
$\hat\xi = 0.038$ (NTHMB sf), $0.118$ (BRIDG sf), $-0.058$ (BRIDG bf) ---
quote these when asked; with $\sim 2250$ exceedances the standard error is
roughly $(1+\xi)/\sqrt{n} \approx 0.025$ before clustering inflation, so all
three are within a couple of standard errors of the exponential boundary
$\xi = 0$: moderately light tails, no evidence of wild heavy-tailedness.
Scale trends: likelihood-ratio statistics $108.2$, $35.4$, $89.3$ on one
degree of freedom --- all far beyond any conventional threshold ($\chi^2_1$
at the 0.001 level is 10.8). NTHMB's mean scale by era: $29.0 \to 63.7$ ---
the typical overshoot more than doubles; BRIDG rises $8.5$ to a
mid-century $13.8$ then eases to $12.0$; baseflow $0.73 \to 2.1 \to 1.7$.

\textbf{The two boxplots, and what each proves.} Raw excesses by month show
strong seasonality --- freshet overshoots are physically larger --- which is
why $\sigma$ must be seasonal. Excesses \emph{divided by} the fitted
$\sigma(t)$ should be month-free if the seasonal scale absorbed it: the boxes
are level, which is the professor's ``excesses no longer seasonal'' check
passing in the body. The outliers are not level --- winter months contain
standardized excesses far larger than summer's --- which is the tail story
below.

\textbf{Question: ``is the constant shape defensible?''} Here is the honest
answer, with its diagnostic. Refit the GP per season on the standardized
excesses: DJF $\hat\xi = 0.271$ ($n = 478$), MAM $0.193$, JJA $-0.180$, SON
$0.051$. Winter is heavier-tailed than summer, and the same signal appears
independently in the seasonal Exp(1) panels (the DJF excesses bow above the
line) and in the standardized-excess outliers. Per-season standard errors are
roughly $0.06$ before clustering inflation, so DJF versus JJA is a real
contrast, not noise --- but freeing $\xi$ by season is a modelling decision
with real cost (short winter records, harder identifiability), and it should
be made with the professor, not slipped in.

\subsection{Station 5: the uniformity gate, in full}
\label{sec:pit}

This is the gate of the whole analysis (v8~\S3, p.~4), it is the check the
professor named, and every follow-up question about it has a precise answer.

\subsubsection{The theorem: why correct margins force a uniform histogram}
\label{sec:pitthm}

\textbf{The probability integral transform.} If a continuous random variable
$X$ has distribution function $F$, then $U = F(X)$ is Uniform$(0,1)$. The
proof is three lines: for $u \in (0,1)$,
\[
\Pr(U \le u) \;=\; \Pr\{F(X) \le u\} \;=\; \Pr\{X \le F^{-1}(u)\}
\;=\; F\{F^{-1}(u)\} \;=\; u,
\]
and ``$\Pr(U \le u) = u$ for every $u$'' is the definition of Uniform$(0,1)$.
The intuition: $F(x)$ asks ``what fraction of this distribution lies below
$x$?'', and a random draw lands below its own 30th percentile exactly 30\% of
the time, below its own 70th percentile exactly 70\% of the time --- at every
level equally. Feeding a variable through its own distribution function
erases the shape of the distribution and leaves pure rank.

\textbf{The non-stationary version, which is our gate.} Every day $t$ owns a
distribution $F(\cdot\,;t)$, and each day's value is pushed through \emph{its
own day's} fitted function: $U_t = F(x_t;t)$. If the fitted margins are
correct, then $U_t$ is Uniform$(0,1)$ for every single day --- a June freshet
day and a February low-flow day become identically distributed --- and a
mixture of uniforms is uniform, so the pooled sample of 47{,}847 values must
look uniform. The contrapositive is the diagnostic: any deviation from
flatness is an accusation that $F(\cdot\,;t)$ was wrong \emph{for some
identifiable set of days}, and \ref{sec:pitband}--\ref{sec:pitmaps} are the
instruments that identify the set.

\textbf{Three fine points that questions tend to probe.}
\begin{enumerate}
\item \emph{Continuity is load-bearing.} If the data carry ties, $F$ has
jumps and $F(X)$ has point masses --- the histogram would spike at the tied
values' probabilities no matter how good the model. This is why the forensics
of Station~\ref{sec:datagate} run first: for streamflow there are no ties;
baseflow's 3\% tie mass is the one place a small spike would be data, not
model.
\item \emph{Uniformity is marginal, not joint.} The theorem makes each day's
$U_t$ uniform; it says nothing about independence, and indeed consecutive
$U_t$ are strongly dependent --- a wet week is a week of high ranks. This
does not move the target (flat at 1); it inflates the \emph{noise} around
the target, quantified next.
\item \emph{The check is slightly in-sample.} The same days that built
$\hat F$ are pushed through it, and a quantile fit calibrates its own rung
shares by construction, so the pooled histogram flatters the fit a little at
the rung levels. The sliced views, the era maps and the rung audit are what
keep the check honest; a fully out-of-sample version (transform each year
through a fit that excluded it) is a possible refinement, not currently
needed.
\end{enumerate}

\subsubsection{The histogram: what a bar is, how flat is flat, and why 40 bins}
\label{sec:pitbars}

\textbf{What a bar is.} With bins of width $w = 1/40 = 0.025$, the bar over a
bin plots $\text{count}/(n\,w)$, so a perfectly uniform sample gives every
bar height 1 \emph{in expectation}. The red line at 1 is that expectation.

\textbf{How much a bar may wobble.} Under \emph{independent} days the bin
count is Binomial$(n, w)$ with $n = 47{,}847$: mean $nw \approx 1196$,
standard deviation $\sqrt{nw(1-w)} \approx 34$, so the bar height has
standard deviation $34/1196 \approx 0.029$ and a $\pm 2$-SD range of about
$\pm 0.06$. But the days are not independent: flow persists, so the $U_t$
move in runs of several days, and clustered counts have variance inflated by
roughly the mean run length. With runs of 4--8 days the standard deviation
multiplies by $\sqrt{4}$--$\sqrt{8} \approx 2$--$3$, giving honest
fluctuation of roughly $\pm 0.15$ per bar. \emph{So ``most bars between 0.85
and 1.15'' is what a correct model looks like at this sample size}; demanding
every bar touch 1 would be demanding less noise than 47{,}000 dependent days
can deliver. The run-length arithmetic is an order-of-magnitude argument ---
which is why the report does not rely on it, but computes the noise band
empirically (next).

\textbf{Why 40 bins, and why not more.} Bar noise scales like
$1/\sqrt{n\,w}$: halve the width (80 bins) and the iid noise grows to
$\pm 0.08$, with clustering to roughly $\pm 0.2$; at 400 bins each bar holds
$\sim 120$ days and the display is pure noise. Go the other way (10 bins) and
the bars are quiet ($\pm 0.03$) but blind to any feature narrower than $0.1$
in probability. Forty is a display-resolution choice, not a law: width
$0.025$ matches the spacing of the ladder's own bottom rungs, so one bar
roughly interrogates one rung interval. And the same theorem is checked at
two other resolutions: the \textbf{QQ plot is the bin-free version} --- it
compares all 47{,}847 order statistics against their reference positions, no
binning at all --- and the \textbf{rung audit is the exact version} at the
thirteen fitted probabilities: the share of days at or below $u(t;\tau_k)$
minus $\tau_k$, at most $0.006$ in absolute value everywhere in our run, with
the only systematic entries at the bottom ($-0.006$, $-0.005$, $-0.005$ at
$\tau = 0.01$: the 1\% rung sits slightly low --- the known difficulty of
extreme quantile levels, and the cause of the short first bar).

\subsubsection{The band figure: a noise interval on every bar}
\label{sec:pitband}

\textbf{The question each blue segment answers:} \emph{how much would this
bar move if history had dealt us a different set of years?}

\begin{Verbatim}[fontsize=\small]
STEP 1  compute the observed 40 bar heights once
STEP 2  repeat 200 times:
          draw 131 year labels with replacement from {1969, ..., 2099}
          keep the days of the drawn years        # whole years: the block
          recompute the 40 bar heights on this replicate
STEP 3  for each bar: the segment spans the 2.5% and 97.5% quantiles
        of its 200 replicate heights
STEP 4  flag (red dot) every bar whose whole segment sits off 1
\end{Verbatim}

\textbf{Why whole years are the resampling unit.} The noise to be measured
comes from serial dependence --- flood clusters, wet seasons --- and a year is
the natural block that keeps a full seasonal cycle and its clusters intact.
Resampling single days would destroy exactly the dependence that creates the
noise and produce a band that is too narrow (the iid $\pm 0.06$), flagging
noise as structure.

\textbf{How to read the figure.} Grey bars: observed heights. Blue segments:
the year-block interval per bar. Red line: the target 1. Red dots: bars
whose interval excludes 1 --- deviations that no plausible re-draw of years
repairs, i.e.\ structural; those are the bars the residual maps then
attribute. Two honesty notes. First, $\hat F$ is not refit inside each
replicate, so the band measures sampling variability of the days, not
uncertainty of the fit --- the right yardstick for ``is this bar's deviation
year-to-year luck.'' Second, the implemented draw keeps a drawn year once
even when it is drawn twice, a subsampling flavour of the bootstrap that
makes the bands slightly \emph{wider} than the exact multiset version --- the
conservative direction for a diagnostic (it flags too few bars, never too
many).

\subsubsection{The residual maps: attributing every off bar}
\label{sec:pitmaps}

\textbf{The construction.} A contingency table: rows are slices (the 12
months, or the 6 eras), columns are the 40 bins, and the cell count $O$ is
``how many January days landed in bin $(0.025, 0.050]$.'' The null
hypothesis --- margins correct --- says each slice spreads \emph{its own}
days evenly over the bins, so the expected count is
$E = (\text{row total})/40$. The colour of a cell is the Pearson standardized
residual
\[
R \;=\; \frac{O - E}{\sqrt{E}},
\]
approximately standard normal under the null (row totals put $E$ near 100
per month cell, near 200 per era cell): white $\approx 0$ is perfect, red is
too many days, blue is too few. $|R| \le 2$ is ordinary noise; because
serial clustering inflates cell noise and 480 cells bring multiplicity (the
largest of 480 null residuals is around 3 by chance alone), the working rule
is that $|R| \ge 4$ is structure --- and the eye should look for
\emph{stripes}, not lone cells.

\textbf{The reading rule --- this is the point of the figure.} A horizontal
stripe along a \emph{month} row means that month is systematically
mis-transformed: residual seasonality, the one outcome the professor's
verdict forbids. A vertical stripe down a \emph{bin} column means one rung is
miscalibrated for all months alike: a ladder-tuning matter, not seasonality.
A stripe along an \emph{era} row means drift the slow basis missed. The map
therefore does not just detect non-uniformity --- it \emph{classifies its
cause}.

\textbf{Our numbers.} No month row shows a stripe --- the seasonal complaint
is answered. The largest cells: BRIDG, era $[1969,1991)$ at the top bin,
$R = +15.2$ --- too many first-era days counted as extreme, the same
first-era threshold miscalibration the local 5\% table shows and the tensor
design repairs; NTHMB, January at bin $(0.025, 0.050]$, $R = +7.1$, a
bottom-rung neighbourhood effect consistent with the low 1\% rung; NTHMB,
era $[2013, 2035)$ at the bottom bin, $R = -7.5$, the same rung read by era.

\textbf{What to say.} ``The PIT is uniform to within the sampling noise that
47{,}000 serially dependent days allow, in each season and each era
separately; the report tests every bar against its own year-block noise band
and attributes every genuine deviation --- and what survives is not
seasonal: it is the slightly-low 1\% rung and the first-era threshold level,
both identified and both with a named repair.''

\subsection{Station 6: the dependence model and the model-free evidence}
\label{sec:depclaim}

\textbf{The model in words.} On the common Laplace scale, write $Y_C$ for the
conditioning station and $Y_X$ for the associated one. Given that $Y_C$ is
extreme ($Y_C = y > v$), the model says
\[
Y_X \;=\; a(t)\,y \;+\; y^{b(t)}\bigl\{m(t) + s(t)\,Z\bigr\},
\qquad Z \sim G
\]
(v8~\S5, eq.~(8), p.~11). Read it as: when one river is extreme at rarity
$y$, the other typically reaches the fraction $a(t)$ of that rarity, with
scatter growing like $y^{b(t)}$, and the leftover randomness $Z$ carries no
formula --- it is represented by the residuals the data themselves produced.
The form is not arbitrary: Heffernan and Tawn showed that for essentially all
standard dependence models, the correct normalisation of $Y_X$ given
$Y_C = y$ large is location $a\,y$ and scale $y^{b}$ --- so the parametric
part is the \emph{shape theory dictates}, and everything data-driven lives in
the four curves and the residual distribution.

\textbf{Why $a$ is the number that matters.} At $a=1, b=0$ the model reads
$Y_X \approx Y_C + Z$: asymptotic dependence --- the link that never fades
however extreme the event. $a=0, b=0$ is exact independence. Everything
strictly between is asymptotic independence with strength $a$; the Gaussian
copula with correlation $\rho$ sits at $a = \rho^2$, $b = \tfrac12$. On the
Laplace scale both tails are exponential, so a single form covers negative
association too ($a \in [-1,1]$; Keef et al.\ 2013). Because $a$ is a
function of $t$, the model can say the \emph{class} of dependence is
changing --- which is the paper's thesis, and Case~4 was selected because the
screening said exactly that.

\textbf{The model-free evidence, before any fitting.} The windowed
$\chi(0.95)$: in each six-year window, count the days on which \emph{both}
stations exceed their own pooled 95th percentile, divide by the days in the
window, divide by 0.05. This estimates
$\Pr\{\text{BRIDG extreme} \mid \text{NTHMB extreme}\}$ from ranks alone:
independence gives 0.05/0.05 $\to$ it hovers near $0.05/0.05 \cdot 0.05
= 0.05$ divided by 0.05, i.e.\ the value $0.05$ shown as the dashed line
after normalisation; perfect dependence gives 1. Ours climbs from $0.253$
(mean of the first ten windows) to $0.403$ (last ten). \emph{The joint
extremes are becoming more frequent, and no model was involved in saying
so.} This is the yardstick every fitted dependence result must face.

\subsection{Station 7: the conditional-simulation display, in full}
\label{sec:pseudo}

This figure carries the most explanation per square centimetre, so it gets
the same four-part treatment as the uniformity gate: what it tests, the
generator with its distribution theory, how to read every visual feature,
and the display decisions. If you can present this section, the functionals
of Station~\ref{sec:funcs} come free, because they are readings off the same
generated cloud.

\subsubsection{What the figure claims to test}

The fitted dependence model is a machine for answering ``given NTHMB is
extreme, what does BRIDG do?'' The pseudosample figure runs that machine
\emph{forward}: it manufactures synthetic days from nothing but the fitted
parameters and the stored residuals, and overlays them on the real days. Two
properties are then visible to the naked eye. \emph{Agreement}: where data
exist, the synthetic cloud should live where the data live --- same centre,
same spread, same tilt. \emph{Extrapolation}: beyond the largest observed
events the cloud keeps going, and that extension is precisely what the model
would answer if asked about a 200-year event --- shown, so its plausibility
can be judged, rather than asserted inside an integral. This is Heffernan and
Tawn's own Figure~3 device, and it is the eye's version of goodness-of-fit
for a conditional model.

\subsubsection{The generator, step by step, with the theory}

\textbf{Step 1: draw the conditioning value.} We need $Y_C$ from its
distribution \emph{given} $Y_C > v$. Two facts make this exact. First, after
the marginal transform, $Y_C$ is standard Laplace, whose upper tail is
$\Pr(Y > y) = \tfrac12 e^{-y}$ for $y > 0$ --- exactly exponential. The
exponential is memoryless, so the excess above any threshold $v \ge 0$ is
itself standard exponential: $\Pr(Y > v + w \mid Y > v) = e^{-w}$. Second,
the inverse-CDF method: if $U$ is Uniform$(0,1)$ then
$\kappa + (1-\kappa)U$ is Uniform$(\kappa, 1)$, and pushing it through the
Laplace quantile function gives
\[
\Pr\bigl[q_{\mathrm{Lap}}\{\kappa + (1-\kappa)U\} \le y\bigr]
= \frac{F_{\mathrm{Lap}}(y) - \kappa}{1-\kappa},
\]
which \emph{is} the conditional CDF of $Y_C$ given $Y_C > v$ (with
$v = q_{\mathrm{Lap}}(\kappa)$, $\kappa = 0.95$, $v = 2.303$). No
approximation: the draw is from the fitted marginal law restricted to the
exceedance region.

\textbf{Step 2: draw the residual --- and why the residual table exists.}
During fitting, for every historical day on which the conditioning station
exceeded $v$, the model equation is solved backwards:
\[
z \;=\; \frac{\bigl(y_X - a(t)\,y_C\bigr)/y_C^{\,b(t)} - m(t)}{s(t)} ,
\]
i.e.\ subtract the fitted level, divide by the fitted growth factor, remove
the residual location and scale. What remains is how far that day sat above
or below its fitted expectation, in standard units, with the effects of
rarity and date both removed. These numbers --- one per historical extreme
day, 2264 of them --- \emph{are} the estimate of $G$. No density is fitted to
them, deliberately: a parametric $G$ would smooth away exactly the tail
irregularities that matter, while the empirical table keeps every skewness
and heavy-tail feature the rivers actually produced. Generation samples one
row with replacement. The licence for reusing a 1975 residual at a 2095 date
is the model's load-bearing assumption --- after the four curves remove the
covariate effect, the residual law no longer depends on $t$ --- and that
assumption is testable and tested: the residual-by-era boxplots must show
equal boxes, and the residual-versus-$y_C$ scatter must show a flat running
mean (both in the report).

\textbf{Step 3: assemble at a chosen date.} Evaluate the four curves at one
day --- reference year $Y$, day-of-year 187, the freshet peak --- and compute
$y_X = a\,y_C + y_C^{\,b}(m + s\,z)$. Every circle is one repetition of
steps 1--3; the panels draw 600.

\textbf{Step 4 (original-scale figure only): push back to flow units.} A
simulated $y_X$ is converted to a probability $p = F_{\mathrm{Lap}}(y_X)$ and
then through the \emph{inverse} of that day's fitted marginal:
below-threshold probabilities invert the ladder interpolation, $p > 0.95$
inverts the GP tail, $p < 0.01$ inverts the exponential lower piece. This is
the same $F(\cdot\,;t)$ machinery as the transform, run backwards, so the
marginal calendar (freshet high, winter low, later years shifted) is restored
exactly.

\subsubsection{The Laplace-scale panel, element by element}

The panel titled ``model pseudosample of $(Y_C, Y_X)$, year 2034.''

\textbf{The axes.} Both variables after the per-day marginal transform, in
standard Laplace units --- the units of \emph{rarity for that day of that
year}. Zero is a median day for its date; $2.30$ is the top 5\% for its
date; $+5$ or $+9$ is far tail. The point of these axes: the calendar is
gone --- a February day and a July day, a 1970 day and a 2095 day, are
directly comparable, because each was judged against its own day's
distribution.

\textbf{The grey cloud (all 47{,}847 days).} A diamond-shaped cloud centred
at $(0,0)$ --- the Laplace shape, exponential in both directions --- tilted
up-to-the-right: positive association surviving the transform, at all
levels, not only extremes. Its symmetric, centred look is itself evidence
that the margins were removed correctly (a wrong transform would displace or
bend it). The upper reach ($\sim +11$) exceeds the lower ($\sim -5$):
extreme high flows go deeper into their tail than droughts do here.

\textbf{The dashed vertical line at $v = 2.303$.} The dependence threshold
--- the Laplace image of the 0.95 level, a \emph{single flat number valid
for every day}, possible only because the transform absorbed the
non-stationarity. The model claims nothing left of this line, which is why
no circles appear there: absence by construction, not failure.

\textbf{The blue circles (600 generated pairs, year 2034).} The year is the
axis midpoint --- the most stable region of the fitted curves; it enters
only through $a, b, m, s$, since these axes carry no calendar. Their
anatomy:
\begin{itemize}
\item \emph{Dense at the line, thinning rightward}: the conditioning value
is threshold plus an Exp(1) draw --- mass decays like $e^{-w}$ beyond $v$
--- matching how the grey points themselves thin to the right.
\item \emph{Their centre rises only gently with $x$}: the gentle rise
\emph{is} $\hat a \approx 0.095$ under the current fit; the vertical width
is $s$ times the residual table's spread, widening slowly because
$\hat b \approx 0.43$.
\item \emph{Circles at low $y$, even below zero}: not an error --- the
finding. Asymptotic independence means BRIDG can be an unremarkable day
while NTHMB is extreme, and the grey data right of the line show the same
low points. A model forcing every circle high would be claiming asymptotic
dependence, which this pair does not have.
\item \emph{The occasional very high circle} (near $(5, 9)$): a heavy
residual drawn from the table --- the empirical $G$ kept the rivers' real
skewed, heavy residuals, which a fitted Gaussian would have erased.
\item \emph{The agreement test, by eye}: compare circles with grey points
in the same $x$-range --- same centre, same spread, same skew. Circles
reaching $x \approx 8$--$9$, where the grey points run out, are the
extrapolation, on display.
\end{itemize}

\textbf{The honest caveat belonging to this panel.} These circles were
generated under the dependence fit that Section~\ref{sec:defect} shows never
left its starting value ($\hat a$ flat at 0.095). The cloud still overlays
the data credibly because $m$, $s$ and the empirical residuals absorb much
of what a mis-estimated $a$ leaves behind --- which is exactly why the eye
test alone cannot settle $a$-versus-$m$, and why the optimiser repair plus
the bootstrap, not this picture, will carry the trend claim.

\subsubsection{The original-scale figure, element by element}

The figure titled ``model pseudosamples on the original scale, day-of-year
187, years 1975 and 2095,'' with both axes logarithmic.

\textbf{Why log axes.} The flows span three orders of magnitude (BRIDG 0.5
to $\sim 300$, NTHMB 10 to $\sim 1500$); linear axes would crush the cloud
into a corner. On log axes equal \emph{ratios} are equal distances --- so a
cloud that looks compact on the right actually spans a wide absolute range.

\textbf{The grey cloud.} Three features worth naming. It rises diagonally:
association at all levels. It is made of fine \emph{parallel threads},
clearest lower-left: recession curves --- consecutive days of one storm or
melt event trace a smooth trajectory as both rivers drain, one thread per
event, made visible by the log scale. And it has \emph{broad branches}
rather than one ellipse: rising-limb, falling-limb and different seasons
occupy slightly different tracks (hysteresis) --- a reminder that the body
of the joint distribution is complicated, which is why the model only
claims the conditional tail.

\textbf{The black triangles (one per year).} For each of the 131 years, the
day NTHMB hit its annual maximum, plotted at that day's (NTHMB, BRIDG)
pair; they pile up at NTHMB $\approx 500$--$1100$, BRIDG $\approx
80$--$250$. They exist for one reason: the meeting's complaint was that
simulated points came out smaller than real extremes, so the real extreme
days are put on the plot and the clouds must land on and extend their
territory --- checkable by eye. (The stray low triangle near $(350, 40)$: a
weak flood year, real, not an error.)

\textbf{The blue circles (year 1975, day 187).} Each is one run of the
generator with both coordinates pushed back through \emph{that day's}
inverse marginals. Two visible properties: the column is \emph{narrow in
$x$} --- the conditioning draw concentrates just above the threshold
(exponential decay of mass) and 1975's day-187 margin maps that
concentration into a tight band starting at the day's 95th percentile
($\sim 650$); and its \emph{vertical spread} ($\sim 80$--$250$) matches the
triangles' spread --- the calibration check passing.

\textbf{The red circles (year 2095, same day).} Everything that differs
from blue is the fitted change across the century: the cloud shifts
\emph{right} (NTHMB's margin intensifies --- threshold slope $+61.8$ per
unit axis, GP scale $29 \to 64$), stretches into territory with no grey
dots --- the extrapolation, displayed --- and rises slightly. Honest note to
volunteer if pressed: with the current stuck fit $a(t)$ is flat, so nearly
all visible blue-to-red movement is the \emph{margins} changing; the
dependence component joins after the optimiser repair.

\textbf{One-sentence summaries.} Laplace panel: ``on the common rarity
scale, the model's synthetic extreme days --- parameters plus recycled
residuals only --- are statistically indistinguishable from the real days
where data exist, and beyond the data they show exactly the extrapolation
the model would use for rarer events.'' Original scale: ``the same synthetic
days placed on the real flow scale at the freshet peak of an early and a
late year: they sit where the real annual-maximum days sit, and the red
cloud is what the fitted century of change does to a co-flood.''

\subsubsection{The display decisions, each with its reason}

\textbf{Why the conditioning day is the freshet peak (day 187).}
Conditioning on a randomly picked date lands mostly in low-flow seasons; the
back-transform then returns modest flows, and the simulated points come out
smaller than the real annual extremes --- precisely the complaint raised in
the meeting. Real annual maxima occur on particular days of the season, so a
display meant to be compared with them must condition on the day of year.
The original-scale figure overlays the observed annual-maximum days as black
triangles so the comparison is checkable by eye: the clouds now land on and
extend the triangles' territory.

\textbf{Why no year colours on the standard-scale panel.} The marginal
transform absorbs the year: after each day passes through its own day's
distribution, all days follow the same Laplace law, so 1973 and 2094 are
\emph{by construction} indistinguishable on those axes, and colouring them
would draw a distinction the axes cannot carry. (The joint drift, if any,
lives in $a(t), b(t)$ --- and is quantified by the fitted curves and the
functionals, not by colours here.)

\textbf{Why the year contrast is legitimate on the original scale.} There the
margins genuinely differ by year --- the same Laplace pair maps to different
flows in 1975 and 2095 --- so blue-versus-red carries real information: the
whole fitted change, margins and dependence together. Which scale the paper
should show is the open display question for the professor; the eventual
headline is RCP~4.5 against RCP~8.5 on this same figure, once that export is
in the pipeline.

\textbf{Why years 1975, 2034 and 2095 --- deliberately not extreme
years.} A display date has two coordinates chosen by different logic. The
\emph{day-of-year} coordinate (187) is chosen to be extreme, on purpose and
by computation, not by hand: it is the day at which the fitted threshold's
climatological average peaks, so the simulation is comparable with real
annual maxima. The \emph{year} coordinate is structural, not extreme:
\code{start\_year + 6}, \code{round(mean(range(years)))} and
\code{max(years) - 4} give 1975, 2034 and 2095 --- early, middle and late
positions fixed by arithmetic before any result is seen. Why those
positions: early-versus-late is the figure's job (same conditioning event,
same day of year, only the year differs, so the blue-to-red difference is
purely the fitted change); a few years \emph{inside} the axis ends because
spline fits are least stable at boundaries; the middle year for the
single-cloud panel because it is the most stable region and, on the
standard scale, any year works almost by construction. The distinction
matters: picking a year \emph{because it contained a famous flood} would be
selection bias --- choosing the display after seeing outcomes; axis
positions declared in advance are illustration (v8~\S6: ``choosing a date
to display is illustration, not selection bias, provided the questions the
figure answers are written down first''). And the conditioning \emph{value}
is never chosen at all --- it is integrated over: $Y_C$ is drawn from the
whole top-5\% region, so the display conditions on ``extreme,'' not on one
historical event.

\textbf{Monte Carlo error, in case it is asked.} With 600 circles the cloud's
visual boundary wobbles; with the $n_{\mathrm{mc}} = 10{,}000$ draws used for
the functionals, a probability near $0.35$ carries a Monte Carlo standard
error of $\sqrt{0.35 \times 0.65/10{,}000} \approx 0.005$ --- negligible
against everything else.

\subsection{Station 8: return levels}
\label{sec:rl}

\textbf{The definition, and why it needs a reference year.} Under a trend
there is no single ``100-year flood.'' Fix a station, a period $T$ and a
reference year $Y$: freeze the fitted curves as they stand in year $Y$, let
$M_Y$ be the annual maximum of a year with those parameters, and define
$x_T(Y)$ by $\Pr\{M_Y > x_T(Y)\} = 1/T$ (v8~\S6, eq.~(10), p.~14). Every
return-level statement is a probability statement about a year with a stated
climate --- the waiting-time reading (``we will wait about $T$ years'') is
false under a trend and is never used.

\textbf{The construction, in three moves.} (i) For $x$ above the threshold on
day $t$, the daily exceedance probability is the chance of being an
exceedance day at all times the GP tail beyond $x$:
$p(x,t) = 0.05 \times \{1 + \xi (x - u(t))/\sigma(t)\}_+^{-1/\xi}$
(v8~eq.~(11)). (ii) Chain the days of year $Y$:
$\Pr(M_Y \le x) = \prod_t \{1 - p(x,t)\}$ --- exact if days were
independent. (iii) They are not: extremes arrive in clusters, so the year
holds fewer independent opportunities than days, and the classical extremal
index result says the dependent maximum behaves like the independent one
raised to the power $\theta \in (0,1]$:
$\Pr(M_Y \le x) \approx [\prod_t \{1-p(x,t)\}]^{\theta}$ (v8~eq.~(13),
p.~15), where $1/\theta$ is the mean cluster length. Ignoring $\theta$
overstates the annual exceedance probability and \emph{inflates} the return
levels --- wrong even though it errs on the cautious side.

\textbf{Our numbers.} $\hat\theta = 0.244$ at NTHMB (mean cluster 4.1 days),
$0.127$ at BRIDG (7.9 days), with our implementation and
\code{extRemes} agreeing to three decimals. The two computational routes ---
root-finding on the formula versus simulating thousands of synthetic years
and reading the empirical quantile --- agree to $0.26\%$, $0.32\%$, $0.98\%$
at $T = 10, 50, 100$ when clustering is switched off, which validates the
chaining machinery before the correction is trusted.

\subsection{Station 9: the three conditional functionals}
\label{sec:funcs}

The professor's three summaries --- the conditional mean
$\mathbb{E}(Y_t \mid X_t \ge x)$, the conditional quantiles, and the band
probability $\Pr(a \le Y_t \le b \mid X_t \ge x)$ --- are all read off one
generated cloud per reference year, produced by exactly the generator of
Station~\ref{sec:pseudo} with $n_{\mathrm{mc}} = 10{,}000$: the mean is
computed \emph{after} inverting to flow units (the mean does not commute with
the nonlinear transform, so it must be taken on the flow scale); the
quantiles are empirical quantiles of the generated flows; $P_{co}$ --- BRIDG
above its own day's threshold while NTHMB is extreme --- needs no inversion
at all, because on the Laplace scale the day's threshold is the same flat
number $y^* = 2.303$ for every day, so it is just the fraction of the cloud
above $y^*$; likewise $P_{band}$ with the band bounds at the day's marginal
0.95 and 0.995 quantiles, which on the Laplace scale are the fixed numbers
$2.303$ and $4.605$. All bounds and levels are declared parameters, because
they were the professor's free choices.

\textbf{Status.} The machinery is verified --- the model-generated
conditional mean matches the plain empirical average over observed
exceedance days to $0.5\%$ (early era) and $8.4\%$ (late era), the HT04
Table-5 discipline --- but the values currently printed
($P_{co} \approx 0.35$, flat in $t$) are computed from the stuck dependence
fit of Section~\ref{sec:defect} and must be re-generated after the repair.
Present the machinery, not the numbers.

\subsection{Station 10: uncertainty}
\label{sec:boot}

Whole years are resampled with replacement, keeping each year's days together
and in order --- the block that preserves within-year clustering --- and
threshold, GP tail, extremal index and dependence model are \emph{refitted}
inside every replicate; the spread of the recomputed quantities is the
uncertainty band. This is different in kind from Station~5's band (years
re-drawn, nothing refitted: a sampling-noise yardstick) and from Route~B
(generation from a fitted model: Monte Carlo, not uncertainty). The knit runs
$B = 25$ replicates for speed; the paper-scale run is $B = 2000$ with every
data-driven selection (smoothing weights, model sizes, the lag) re-selected
inside each replicate --- selection outside the loop makes bands too narrow,
because it treats a chosen quantity as if it had been known in advance.

\section*{Part II --- Inside the code}
\addcontentsline{toc}{section}{Part II --- Inside the code}

Every algorithm: what it computes, the pseudocode, why that is the right
computation, and what verifies it. Where a published package does the work,
what it does internally is explained --- ``we use qgam'' is not an
explanation.

\subsection{The one formula behind every parameter curve}

Every parameter --- threshold, scale, the four dependence curves --- is one
formula with different ingredients:
\[
\eta(t) = \sum_{k=1}^{p} \gamma_k B_k(t),
\]
where the $B_k$ are fixed, known basis functions chosen before fitting and
the $\gamma_k$ are the only estimated quantities (v8~\S4.1, eq.~(7), p.~5).
One column: a constant (the stationary model). Two columns $[1, t]$: a
straight line. Cubic B-spline bumps at knots along the axis: the flexible
curve. The bases are nested --- constant $\subset$ line $\subset$ spline ---
which is what makes model-size comparison meaningful. A cubic B-spline bump
is nonzero over only four knot intervals, rises to one peak, and returns to
zero; at any date exactly four bumps are nonzero and they sum to one, so the
curve at a date is a local weighted average of four coefficients, each
estimated from the thousands of days under its own bump. Neighbours replace
repeats; smoothness is the assumption that makes that legitimate. The
seasonal component adds cyclic basis functions in day-of-year (period one
year, derivatives matching across the year boundary); the tensor version
multiplies the two bases into a surface.

\subsection{The quantile ladder in code: qgam, line by line}
\label{sec:qgamcode}

The target is unchanged from Station~\ref{sec:threshold}: for each level
$\tau$ in $\{0.01, 0.025, 0.05, 0.1, \ldots, 0.9, 0.95\}$, the curve
minimising pinball loss plus roughness --- and minimising pinball loss
\emph{is} quantile estimation, by the derivative identity proved there.
What follows is every qgam-touching line, in the code-understanding format.

\textbf{Entry 1: the model formula.}
\begin{Verbatim}[fontsize=\small]
fm <- list(y ~ s(doy, bs = "cc", k = 12) + s(tt, bs = "cr", k = 25),
           ~ s(doy, bs = "cc", k = 8))
\end{Verbatim}
\code{s()} declares a smooth term in mgcv's language. \code{bs = "cc"} is
the \emph{cyclic} cubic spline: the basis is built so the curve and its
derivatives match across the year boundary --- December 31 joins January 1
with no seam; one seasonal shape, learned from all 131 years at once.
\code{bs = "cr"} is the ordinary cubic spline for the trend; $k = 25$ over
131 years is a knot roughly every 5.5 years, matching the project's
spacing decision. \textbf{What $k$ is and is not:} $k$ is the basis
dimension --- a \emph{ceiling} on flexibility, not the flexibility itself;
the actual wiggliness is chosen by the smoothing penalty (Entry 3), so a
generous $k$ with automatic smoothing is the standard idiom. \textbf{The
second formula is not a second quantile model}: it tells qgam how the
conditional \emph{spread} of $y$ varies, because the smoothed loss has a
per-observation bandwidth that must scale with local spread (Entry 3, step
1). Flow variance is massively seasonal --- freshet days scatter over
hundreds of units, winter days over a few --- so the variance model is
cyclic in day-of-year. Omitting it is qgam's classic failure mode: a
constant bandwidth gives the fitted 0.95 curve season-dependent bias
(their own motorcycle-data figure demonstrates it).

\textbf{Entry 2: the pass-through knots.}
\begin{Verbatim}[fontsize=\small]
argGam = list(knots = list(doy = c(0.5, 366.5)))
\end{Verbatim}
Everything in \code{argGam} goes to the internal \code{mgcv::gam} calls.
This pins the cyclic boundary knots at half-days: the period is exactly 366
days with the join \emph{between} Dec 31 and Jan 1. Without it, mgcv sets
the cycle boundary at the observed data range (1 to 366), which would
identify day 1 \emph{with} day 366 --- gluing New Year onto New Year's Eve
as the same covariate value instead of neighbours.

\textbf{Entry 3: the fit itself.}
\begin{Verbatim}[fontsize=\small]
fitM <- mqgam(fm, data = d, qu = taus, argGam = kn)
\end{Verbatim}
\code{mqgam} fits the model at all 13 levels, storing the shared components
once (why it beats 13 \code{qgam} calls). Internally, in order:
\begin{enumerate}
\item \emph{Preliminary location--scale fit.} A Gaussian \code{gaulss} GAM
of mean and variance (the variance from our second formula), plus a
flexible density (sinh--arcsinh) for its standardized residuals. Purpose:
the pinball loss has a kink, which blocks smooth optimisation; qgam
replaces it with the smoothed \emph{ELF loss}, and the amount of smoothing
is \emph{computed, not tuned} --- a formula minimising the asymptotic mean
squared error of the coefficients, evaluated with this density estimate,
giving a \emph{per-observation} bandwidth: wide in the freshet, narrow in
winter. Their theoretical result: the optimally smoothed loss is
statistically \emph{better} than the raw pinball, not merely convenient.
\item \emph{Learning-rate calibration, per level} --- the expensive step
and the printed dots (``\code{Estimating learning rate... qu = 0.5
......done}''). The rate $1/\sigma_0$ balances loss against the smoothing
prior; qgam selects it by Brent search where \emph{every dot is a complete
penalized fit} at a trial $\sigma_0$, scored so that the intervals for the
quantile curve achieve roughly their advertised coverage. This is why 13
rungs $\times$ 5 series was slow: calibration, not fitting.
\item \emph{Smoothing selection.} Given $\sigma_0$ and the bandwidth, the
roughness weights $\lambda$ (one per \code{s()} term) are chosen by
Laplace-approximate marginal likelihood --- the step that replaced our
cross-validation loop at this stage (the head-to-head is
Section~\ref{sec:qgamvscv}).
\item \emph{Final fit} at the chosen $\sigma_0, \lambda$ by penalized
iteratively re-weighted least squares --- mgcv's standard engine.
\end{enumerate}
Recorded costs, all handled by our code: no non-crossing constraint
(Entry 5); calibration theory assumes independent data, so with 4--8-day
clusters qgam's intervals are a working approximation and \emph{all} our
uncertainty comes from the year-block bootstrap; and extreme-$\tau$ bias
--- our rung audit shows exactly this cost, the $\tau = 0.01$ rung sitting
$\sim 0.5\%$ low.

\textbf{Entry 4: extracting the ladder.}
\begin{Verbatim}[fontsize=\small]
U <- sapply(taus, function(q) as.numeric(qdo(fitM, q, predict, newdata = d)))
\end{Verbatim}
\code{mqgam} stores its 13 fits stripped of shared components to save
memory, so they cannot be used directly; \code{qdo} re-attaches the shared
pieces to the level-$q$ fit and applies any standard gam function --- here
\code{predict} on all 47{,}847 days. \code{sapply} stacks the 13 curves
into the $N \times 13$ matrix $U$. Downstream, \emph{only $U$ is used} ---
which is what makes the qgam route and the fallback interchangeable.

\textbf{Entry 5: the non-crossing audit and repair.}
\begin{Verbatim}[fontsize=\small]
ncross <- sum(apply(U, 1, function(r) any(diff(r) <= 0)))
if (ncross > 0) { Us <- t(apply(U, 1, sort))
                  msize <- max(abs(Us - U)) / sd(x); U <- Us }
\end{Verbatim}
Per day, check the 13 rung values are strictly increasing; count the days
with any violation; repair by sorting each day's rungs into order ---
monotone rearrangement, which provably never increases estimation error
(Chernozhukov, Fern\'andez-Val, Galichon 2010). \code{msize} records the
largest change the sort made, in sd units --- the number separating
``hairline overlaps'' from ``real fitting problem.'' Our run: 1147 / 1520 /
1481 crossing days, largest repair 0.013 / 0.006 / 0.009 sd --- cosmetic.
The returned \code{lam\_u = NA} marks that qgam kept its smoothing
internal; downstream code checks for the NA.

\textbf{Entry 6: the shell --- fallback and cache.} The \code{tryCatch}
around the body degrades to the project \code{quantreg} fitter on any
error, announced with \code{cat("NOTE: ...")} --- \code{cat}, not
\code{message()}, because chunk options suppress messages and a silent
fallback once hid a wrong route. The cache wrapper saves the result to
\code{cache\_ladder\_<tag>\_<route>\_<N>\_<sum(x)>.rds}: the name carries
the route and a data fingerprint, so changing the data, axis or route
changes the filename and forces a refit automatically;
\code{refit\_ladders: true} forces one manually; the knit prints which
cache file it used --- provenance on the record.

\textbf{Entry 7: the fast-calibration variant.}
\begin{Verbatim}[fontsize=\small]
cal <- tuneLearnFast(fm, data = d, qu = c(0.05, 0.50, 0.95), argGam = kn)
lsig_of <- approxfun(c(0.05, 0.50, 0.95), cal$lsig, rule = 2)
U <- sapply(taus, function(q) as.numeric(predict(
       qgam(fm, data = d, qu = q, lsig = lsig_of(q), argGam = kn),
       newdata = d)))
\end{Verbatim}
Entry 3's step 2 runs per rung and dominates the cost.
\code{tuneLearnFast} runs \emph{only} that calibration, at three anchor
levels; \code{approxfun} interpolates the selected $\log\sigma_0$ linearly
in $\tau$ (the learning rate varies slowly and smoothly in $\tau$, which is
what makes interpolation safe); passing \code{lsig} to \code{qgam} skips
calibration entirely, so each rung is one plain fit. Net: 3 calibrations
$+$ 13 fits instead of 13 calibrations --- a 3--4$\times$ cut, same model.

\textbf{Entry 8: the tensor experiment.}
\begin{Verbatim}[fontsize=\small]
fte <- qgam(list(y ~ te(doy, tt, bs = c("cc","cr"), k = c(10, 15)),
                 ~ s(doy, bs = "cc", k = 8)),
            data = d, qu = 0.95, argGam = kn)
\end{Verbatim}
\code{te()} replaces ``seasonal shape $+$ trend'' with a smooth
\emph{surface} in (day-of-year, year): the marginal bases (10 cyclic
$\times$ 15 trend) are multiplied into 150 surface functions with
\emph{separate} penalties in each direction --- two smoothing parameters,
both by Entry 3's machinery --- cyclic in the doy direction only. The
additive model is a special case of the surface, so nothing can be lost,
and the smoothing shrinks back toward additivity unless the data pay for
the interaction. Fitted at the single working level $\tau = 0.95$ (one
calibration), cached under its own filename. This is the fit whose season
$\times$ era table turned the first-era misses (0.149, 0.123) into 0.039
--- the evidence for adopting it.

\subsection{qgam versus blocked cross-validation: one problem, two philosophies}
\label{sec:qgamvscv}

Both routes estimate the same curve under the same loss; they differ in the
criterion for one nuisance number, the roughness weight. The big picture in
one sentence: \emph{cross-validation treats smoothing selection as a
prediction experiment --- physically hold data out, measure pinball loss on
it, keep the $\lambda$ that predicts best; qgam treats it as Bayesian
inference --- translate the penalized problem into likelihood-plus-prior,
then keep the $\lambda$ under which the observed data are most probable.}
Everything qgam builds (the smoothed loss, the learning rate) exists to
make that translation legitimate, because a raw loss is not a likelihood
and Bayesian machinery may not touch it directly. Step by step:

\textbf{Step 0, the target: identical.} The $\tau$-quantile curve, defined
by the pinball identity. Nothing about the route changes what is estimated.

\textbf{Step 1, the objective.} CV route: raw pinball $+$
$\lambda\,\times$ squared differences of neighbouring coefficients. qgam:
the same with the kink rounded off (ELF), the rounding set by the
MSE-minimising formula with a per-observation bandwidth --- not a new knob,
and provably an upgrade, not a compromise.

\textbf{Step 2, the inner fit at fixed $\lambda$.} CV route: a linear
program (\code{rq.fit.fnb} on the augmented design, penalty rows appended
as pseudo-data). qgam: with a smooth loss, $\exp(-\mathrm{loss})$ behaves
like a density, so the fit is penalized likelihood via iteratively
re-weighted least squares. Same role, different solver, enabled by step 1.

\textbf{Step 3, choosing $\lambda$ --- the head-to-head.} CV route: a grid
of $\lambda$ times five year-block folds; fit on four fifths, score
held-out pinball, average, argmin. Direct and model-free, but expensive,
one $\lambda$ only (a two-dimensional grid times folds is unaffordable ---
which is why our CV shares a single $\lambda$ across curves), and it stops
at the grid edge (we hit that edge twice, at $10^4$ and $10^6$). qgam: the
penalty is reinterpreted as a Gaussian \emph{prior} --- ``wiggly curves are
improbable a priori,'' with $\lambda$ its precision; integrate the
coefficients out and what remains is the \emph{marginal likelihood of
$\lambda$ itself}: how probable the observed data are under each candidate
smoothness, averaged over every curve that smoothness allows.
Laplace-approximate the integral (LAML), maximise by Newton. No grid, no
folds, no edges, and \emph{several} $\lambda$'s at once (seasonal, trend,
the tensor's two). The bridge to what you know: maximising marginal
likelihood is asymptotically close to optimising leave-out predictive
performance --- the known REML--CV connection --- it reaches the same place
by probability calculus instead of data splitting.

\textbf{Step 4, the step CV does not have: the learning rate.} A true
likelihood fixes the relative weight of data and prior by itself; a
loss-based pseudo-likelihood has no natural scale ---
$\exp(-\mathrm{loss}/\sigma_0)$ is equally valid for any $\sigma_0$, and
$\sigma_0$ controls wiggliness and interval width. This is the price of the
loss-to-likelihood translation (the Bissiri--Holmes--Walker belief-updating
framework is its licence). qgam pins $\sigma_0$ by coverage calibration
(the Brent search, the printed dots). Our CV route never needed this step
because it never claimed intervals --- uncertainty was always the
bootstrap's job; accordingly, the calibrated intervals are the one qgam
output we deliberately do not use.

\textbf{Step 5, serial dependence.} Our blocked folds exist \emph{because}
of dependence: random folds leak a held-out day's neighbours into training
and under-smooth. LAML has no folds, so no leakage of that form --- but it
is not immune either: it counts every day as an independent contribution,
overstating the effective sample size, with the same qualitative risk and
definitely too-narrow intervals. \emph{Neither criterion handles our
dependence honestly by itself}; on both routes the final authority is
out-of-model --- the season $\times$ era tables --- and the year-block
bootstrap for uncertainty.

\textbf{Step 6, what stays ours on either route.} The rung grid, the basis
design, the non-crossing audit and repair, the assembly into $F(x;t)$, and
the verification ladder --- including the comparison of the two routes
themselves.

\begin{center}
\small
\begin{tabular}{@{}p{3.4cm}p{5.3cm}p{5.6cm}@{}}
\toprule
\textbf{Step} & \textbf{Our CV route} & \textbf{qgam's replacement} \\
\midrule
Loss & raw pinball & smoothed pinball (ELF); bandwidth by MSE formula
$+$ variance model \\
Inner fit & linear program (\code{rq}, augmented design) & penalized IRLS
(loss now differentiable) \\
Penalty's meaning & roughness term added to the loss & Gaussian
\emph{prior} on coefficients \\
$\lambda$ selection & grid $\times$ year-block folds, held-out pinball &
marginal likelihood (LAML), Newton, several $\lambda$'s at once \\
Data-vs-penalty weight & implicit, absorbed in $\lambda$ & explicit
learning rate $\sigma_0$, set by coverage calibration \\
Dependence & blocked folds & nothing built in --- bootstrap and
out-of-model tables cover it \\
Uncertainty & block bootstrap (external) & calibrated intervals (built in;
deliberately unused here) \\
\bottomrule
\end{tabular}
\end{center}

The honest closing symmetry, if asked ``which is right'': neither, by
authority. Where the two routes disagree --- top-rung mean gap 0.166 sd ---
the referee is not preference but the out-of-model calibration tables,
where the qgam route with the tensor design currently wins.

\subsection{Blocked cross-validation --- where it still runs, and why blocked}

qgam and evgam select their own smoothing; cross-validation remains only
where no package covers our criterion (our GP fallback fitter; the
dependence stage). It selects one number, $\lambda$:

\begin{Verbatim}[fontsize=\small]
STEP 1  assign whole YEARS to K = 5 folds, round robin (1969->1, 1970->2, ...)
STEP 2  for each candidate lambda in a grid:
          for each fold k: fit on the other folds, score the pinball or
          negative log-likelihood on fold k alone; average the K scores
STEP 3  keep the lambda with the smallest average held-out score
\end{Verbatim}

Whole years, because flood days come in 4--8-day runs: with random-day folds
the neighbours of a held-out day sit in the training set carrying nearly the
same information, the held-out score is optimistic for wiggly fits, and CV
systematically under-smooths. Blocking puts an entire cluster on one side of
every split. Limitation, stated up front: CV answers ``what predicts best''
and can never certify a trend --- the penalty's null space (constants under a
first-difference penalty) is invisible to it --- so significance belongs to
move one, never to $\lambda$.

\subsection{Model size: the two moves}

\textbf{Move one, constant versus linear} --- the only step that can say ``a
trend exists.'' The test matches the stage's criterion. GP tail: a genuine
likelihood, so the likelihood-ratio statistic
$\Lambda = 2(\hat\ell_{\text{lin}} - \hat\ell_{\text{const}})$ is compared to
$\chi^2_1$ (Wilks: under the null, twice the log-likelihood gain from one
extra free parameter is $\chi^2_1$). Threshold: no likelihood exists, so the
slope is judged against its block-bootstrap standard error. Dependence: the
likelihood is a working device, so its $\Lambda$ is only a guide and the
honest test is the bootstrap interval for the slope. \textbf{Move two, linear
versus spline}: no separate test --- the penalised spline contains the line,
so if the fitted spline collapses to a straight shape, the line is reported
and the spline kept as robustness.

\subsection{The GP tail fit}

\begin{Verbatim}[fontsize=\small]
INPUT   excesses (x - u(t)) on the ~2250 days above the threshold
MODEL   log sigma(t) = cyclic seasonal + trend spline   (log => sigma > 0)
        xi = one constant per series
FIT     maximise sum of GP log-densities + roughness penalty on the trend
\end{Verbatim}
Production: \code{evgam}, \code{family = "gpd"}, REML smoothing. Our own
penalised-likelihood fitter runs alongside for the constant-vs-linear LRT and
as the consistency check: relative gaps 0.185 / 0.363 / 1.010 --- the large
baseflow value is two smoothing philosophies disagreeing where $\sigma$ is
tiny (a relative measure inflates), reported not hidden.

\subsection{From two pieces to one distribution; the Laplace map; the inverse}

$F(x;t)$ is assembled per day (v8~\S3, eq.~(5), p.~3): linear interpolation
between rung values in the body; the GP tail above the top rung (continuous
by construction, since $F = 0.95$ at the threshold on every day); an
exponential lower piece below the bottom rung, slope-matched so the lower
tail decays instead of clipping (the old hard clamp at $10^{-5}$ collapsed
the lowest days onto one point at $-2.44$ and drew a flat-run-and-jump QQ
artifact; the instrumented rebuild prints 0 days at the numerical floor).
Then the Laplace map $y = \log(2F)$ for $F < \tfrac12$,
$-\log\{2(1-F)\}$ otherwise (v8~eq.~(6), p.~4). Why Laplace: both tails
exactly exponential and symmetric, so one parametric dependence form covers
positive and negative association ($a \in [-1,1]$), where the Gumbel scale
of the 2004 paper covers positive only. The inverse $q_{\mathrm{inv}}(p;t)$
runs the same assembly backwards (ladder inverse in the body, GP quantile
above, exponential inverse below) and is what returns simulated values to
flow units.

\subsection{The dependence fitter, the Keef constraint, and the defect}
\label{sec:defect}

\textbf{The objective.} Writing $\mu_i = a(t_i) y_{C,i} + m(t_i)
y_{C,i}^{b(t_i)}$ and $\sigma_i = s(t_i)\, y_{C,i}^{b(t_i)}$, minimise over
the exceedance days
\[
\sum_i \Bigl[\log \sigma_i + \tfrac12\{(y_{X,i} - \mu_i)/\sigma_i\}^2\Bigr]
\;+\; \text{roughness penalties},
\]
the negative log of a working Gaussian likelihood. The Gaussian is a scoring
device, not a belief --- Keef et al.\ tested Normal against Laplace working
residuals and kept the Normal --- and the residual is stored as a table, never
believed normal. Links: $a = 2\,\mathrm{plogis}(\eta_a) - 1 \in (-1,1)$;
$b < 1$ enforced; $s = e^{\eta_s} > 0$.

\textbf{The Keef constraint (v8~\S5, eq.~(9), p.~13).} An unconstrained fit
can extrapolate conditional quantiles \emph{above what perfect positive
dependence would give} --- a self-contradiction. The ordering constraint
forbids it, checked at a level $v_K$ above the largest observed conditioning
value, because extrapolation is what needs policing.

\textbf{The defect, precisely.} The constraint was implemented as a hard
wall: objective $= 10^{10}$ outside the feasible set. BFGS estimates
gradients by finite differences; when trial steps land on the wall the
gradient is garbage and the search stops at its starting point. Evidence:
all three model sizes returned the identical objective 2698.6; the trend
statistic was exactly 0.00 with $\hat a$ flat at 0.095; switching the wall
off dropped the objective to 2022.2 (a 676-unit gap --- not a constraint
biting mildly) and produced $\hat a(t): 0.396 \to 0.490$, rising. So the
reported constrained fit is the starting value, not an optimum. One
corollary: \emph{every} constrained linear/spline number in this knit is
suspect --- including the freshet-only linear endpoints ($0.443 \to -0.245$),
fitted through the same wall; the freshet number that is reliable is the
constant $\hat a = 0.227$.

\textbf{What the answer probably looks like.} The unconstrained trend
($0.40 \to 0.49$) and the model-free windowed $\chi$ ($0.253 \to 0.403$)
agree in level and direction. And the constant fit --- which was seeded
sensibly --- agrees with texmex's constrained reference ($a$: 0.095 vs
0.081; $b$: 0.425 vs 0.402), so the constraint machinery is sound in
principle; only the wall is unusable by a gradient method.

\textbf{The repair.}
\begin{Verbatim}[fontsize=\small]
violation(a,b) = size of the ordering breach (0 when feasible)
objective      = working criterion + penalties + kappa * violation^2
fit            = Nelder-Mead first (derivative-free), BFGS to polish,
                 from several starts incl. the unconstrained solution
report         = the violation at the solution (zero = feasible, printed)
\end{Verbatim}

\textbf{The second finding, independent of the defect.} Constant fits per
season of the conditioning day: JJA $\hat a = 0.172$, MAM $0.141$, SON
$-0.011$, DJF $-0.055$. The association lives in the freshet; winter
``extremes'' (extreme \emph{for the season}) carry none. Because the seasonal
threshold spreads conditioning days across all seasons ($537$--$599$ per
season), a pooled $\hat a$ mixes these regimes and is diluted toward zero ---
the freshet-only constant is $0.227$ against the pooled $0.095$. Whether the
science question is season-relative extremes or flood-season extremes is a
decision for the professor; the numbers show it changes the answer.

\subsection{The verification ladder, in one place}

The answer to ``how do I know the code is right'' --- six independent checks,
all printed in the report:
\begin{enumerate}
\item qgam ladder vs our penalised \code{quantreg} ladder: 0.166 sd.
\item evgam tail vs our penalised likelihood: 0.185 / 0.363 / 1.010.
\item Our all-constant dependence fit vs \code{texmex} on its own defaults
(Laplace margins, Keef constraints), both given the identical transformed
pair so the margins cancel: $a$ 0.095 vs 0.081, $b$ 0.425 vs 0.402.
\item With $b$ frozen and curves linear, our fitter vs
\code{mgcv::gam(gaulss)} on literally the same model: conditional means agree
to $3.4 \times 10^{-5}$ of the data range.
\item Route A root-finding vs Route B simulation: $\le 1\%$.
\item The known-answer case: the Case 1 anchor pair must return $\hat a$
near 1; it returns 1.000.
\end{enumerate}
Plus \code{extRemes} $\equiv$ our Ferro--Segers to three decimals, and the
HT04 Table-5 check: model-generated vs empirical conditional mean, 0.5\% and
8.4\% by era.

\section*{Part III --- Question drill}
\addcontentsline{toc}{section}{Part III --- Question drill}

\textbf{``Why quantile regression for the threshold?''} A fixed level is a
different rarity in different seasons and decades, so its exceedances are not
comparable events; fixing the probability and moving the level makes
``extreme'' mean rarity 0.05 on every date --- and minimising the pinball
loss \emph{is} quantile estimation, by the derivative identity
$\frac{d}{du}\mathbb{E}\rho_\tau = F(u) - \tau$.

\textbf{``How is the smoothing chosen?''} Automatically: marginal likelihood
with a calibrated learning rate at the threshold (qgam), REML at the tail
(evgam), blocked cross-validation where no package covers the criterion. The
hand-set quantities are declared design choices: the rung grid, knot
spacing, basis sizes.

\textbf{``Why whole-year folds and blocks?''} Days cluster in 4--8-day runs;
random folds leak a held-out day's neighbours into training, bias the score,
and under-smooth. Years are the natural block containing whole clusters and
whole seasonal cycles.

\textbf{``Why should $F(x_t;t)$ be uniform?''} The probability integral
transform: a continuous variable through its own CDF lands below its
$u$-quantile with probability exactly $u$ --- the definition of uniformity.
Each day uses its own day's fitted CDF, and a mixture of uniforms is
uniform (Section~\ref{sec:pitthm}).

\textbf{``Why 40 bars / why are bars not all at 1?''} Resolution choice;
bar noise $\propto 1/\sqrt{n \times \mathrm{width}}$: $\pm 0.06$ iid,
$\approx \pm 0.15$ with clustering. Bars are tested against year-block
resampling intervals, not against the eye
(Sections~\ref{sec:pitbars}--\ref{sec:pitband}).

\textbf{``What is the heatmap?''} Pearson residuals $(O-E)/\sqrt{E}$ of a
slice-by-bin table; month stripe = seasonality, column stripe = a rung, era
stripe = drift. Ours shows no month stripe (Section~\ref{sec:pitmaps}).

\textbf{``What exactly is $a$?''} The fraction of the conditioning station's
rarity that the other station typically attains on a common standardised
scale: 1 asymptotic dependence, 0 independence, Gaussian copula at $\rho^2$;
time-varying $a$ is the paper's question (Section~\ref{sec:depclaim}).

\textbf{``Where do the simulated points come from?''} An exact draw from the
fitted marginal restricted to the exceedance region (Laplace tail is
memoryless), a residual recycled from the historical table, assembled through
the fitted curves at a chosen date, optionally inverted to flow units ---
Section~\ref{sec:pseudo}, four steps, each with its theory.

\textbf{``Why do the circles only appear right of the dashed line?''} The
model is defined only for conditioning values above $v$; the left region is
not modelled, so no circles belong there.

\textbf{``What would a bad fit look like in the pseudosample figure?''}
Circles systematically above or below the data cloud (wrong $a$ or $m$),
circles fanning wider or narrower than the data (wrong $b$ or $s$), or a
symmetric cloud against visibly skewed data (a parametric $G$ would do this;
the empirical table is why ours does not).

\textbf{``Why condition on day-of-year 187?''} Because annual extremes occur
on freshet days; conditioning on a random date compares simulated low-season
days against observed annual maxima and loses --- the complaint from the
meeting, fixed by design and checkable via the overlaid triangles.

\textbf{``Why years 1975, 2034 and 2095?''} The day of year is chosen to
be extreme --- the climatological freshet peak, computed from the fitted
threshold --- so the simulation is comparable with real annual maxima; the
years are early, middle and late axis positions fixed by arithmetic in
advance. Representative dates, not selected events: picking a year for its
famous flood would be selection bias, and changing these is a one-line
parameter.

\textbf{``CV or qgam --- which is right?''} Both estimate the same curve
under the same loss; they differ only in the criterion for one nuisance
number (Section~\ref{sec:qgamvscv}). Where they disagree --- 0.166 sd at
the top rung --- the referee is the out-of-model calibration table, where
the tensor-qgam route wins.

\textbf{``Is $\xi$ constant defensible?''} A decision, with the diagnostic on
the table: per-season $\hat\xi$ = DJF 0.271, MAM 0.193, JJA $-0.180$, SON
0.051 --- winter heavier, consistent across three independent views, worth a
deliberate discussion rather than a silent change.

\textbf{``Why is the dependence trend not significant?''} It is not yet
tested: the constrained optimiser never left its starting point (identical
objectives, exactly zero statistic), so the flat $\hat a(t)$ measures the
optimiser. The model-free $\chi$ rises 0.25 $\to$ 0.40 and the unconstrained
fit rises 0.40 $\to$ 0.49; the constrained refit with a smooth penalty is the
outstanding fix.

\textbf{``What is the strongest single sentence the marginal stage
supports?''} ``At every series the high quantile and the overshoot scale are
rising significantly, the threshold is calibrated locally by season and era,
and the transform passes its uniformity gate in every slice --- so the
marginal intensification result stands on checks, not assumptions.''

\section*{Part IV --- The fix list}
\addcontentsline{toc}{section}{Part IV --- The fix list}

\begin{enumerate}
\item \textbf{Dependence optimiser (blocking).} Smooth penalty in place of
the wall; Nelder-Mead then BFGS; multiple starts including the unconstrained
solution; print the violation at the solution. Until then, no
dependence-stage number is quotable --- including the functionals.
\item \textbf{Adopt the tensor threshold.} The local table settles it:
first-era misses 0.149 / 0.123 become 0.039 / 0.039.
\item \textbf{Decide the conditioning question} with the professor:
season-relative extremes (current) or flood-season extremes (freshet-only).
The per-season table shows the choice changes the answer.
\item \textbf{Lift the $\tau = 0.01$ rung} (or re-tune the bottom of the
ladder) so the below-rung share matches its nominal 1\%.
\item \textbf{Raise the bootstrap} to $B = 2000$ with selections inside the
loop, once the dependence fit is trustworthy.
\item \textbf{Bring in RCP 8.5}, turning the year contrast into the scenario
contrast the display was rebuilt for.
\end{enumerate}

\end{document}
